\section{Phase 0: Canonical Problem Setting for BLOM}
\label{sec:blom_phase0}

This section fixes the minimal canonical setting used in the first stage of the BLOM theory.
The goal is not to cover the most general formulation, but to define a mathematically stable starting point on top of which the later existence, uniqueness, analytic construction, and convergence results can be built.

\subsection{Problem Statement}
\label{subsec:phase0_problem}

We consider a scalar flat output trajectory
\[
p:[0,\mathcal T] \to \mathbb R,
\qquad
\mathcal T = \sum_{i=1}^{M} T_i,
\qquad T_i > 0,
\]
composed of $M$ polynomial pieces
\[
p_i:[0,T_i] \to \mathbb R,
\qquad i=1,\dots,M,
\]
where each $p_i$ is a polynomial of degree at most $2s-1$.
In Phase~0 we fix
\[
s = 4,
\]
so each segment is a polynomial of degree at most $7$.

Let the waypoint sequence be
\[
q_0,q_1,\dots,q_M \in \mathbb R,
\]
and let the time vector be
\[
T=(T_1,\dots,T_M)^\top \in \mathbb R_{>0}^{M}.
\]
We only prescribe zeroth-order intermediate conditions, i.e.
\[
d_i = 1,
\qquad i=1,\dots,M-1,
\]
which means that the trajectory is required to pass through the intermediate waypoints, but no higher-order derivative is fixed there.

The global multistage minimum-snap reference problem is
\begin{equation}
\label{eq:phase0_global_minco}
\begin{aligned}
\min_{\{p_i\}_{i=1}^{M}} \quad &
\sum_{i=1}^{M} \int_0^{T_i} \left| p_i^{(4)}(t) \right|^2 \, dt \\
\text{s.t.} \quad
& p_i(0)=q_{i-1}, \qquad p_i(T_i)=q_i, \qquad i=1,\dots,M, \\
& p_i^{(\ell)}(T_i)=p_{i+1}^{(\ell)}(0), \qquad \ell=0,1,\dots,2s-2=6, \\
& \hspace{5.1cm} i=1,\dots,M-1.
\end{aligned}
\end{equation}
Equation~\eqref{eq:phase0_global_minco} is not yet the BLOM construction itself; it is the global MINCO-type reference problem against which local constructions will later be compared.

\subsection{Canonical BLOM Window Setting}
\label{subsec:phase0_canonical_window}

In Phase~0 we fix the BLOM window width to
\[
k=2,
\]
namely the \emph{3-point stencil} setting.
For each segment index $i$, define the local window
\begin{equation}
\label{eq:phase0_window}
W(i,2)
=
\{i-1,i,i+1\}\cap \{1,2,\dots,M\}.
\end{equation}
Thus the coefficient of the $i$-th segment will later be defined only from local waypoint--time data associated with at most three neighboring segments.
This is the smallest nontrivial setting that can simultaneously:
\begin{enumerate}
  \item encode local coupling,
  \item exhibit local support explicitly,
  \item remain analytically tractable, and
  \item admit direct numerical verification.
\end{enumerate}

\subsection{Boundary Convention}
\label{subsec:phase0_boundary}

In Phase~0 we adopt only \emph{natural boundary conditions} for local window problems.
For a local window with left boundary piece $p_{\mathrm{left}}$ and right boundary piece $p_{\mathrm{right}}$, we impose
\begin{equation}
\label{eq:phase0_natural_bc}
 p_{\mathrm{left}}^{(\ell)}(0)=0,
 \qquad
 p_{\mathrm{right}}^{(\ell)}(T)=0,
 \qquad \ell=s,s+1,\dots,2s-2,
\end{equation}
with $s=4$, i.e. for $\ell=4,5,6$.
No periodic boundary conditions are considered in Phase~0.
The reason is methodological: before discussing alternative boundary mechanisms, we first need a single canonical definition under which well-posedness and local analytic structure can be studied.

\subsection{Canonical Research Goal of Phase~0}
\label{subsec:phase0_goal}

The sole purpose of Phase~0 is to freeze the mathematical object.
More precisely, the Phase~0 output is the following canonical BLOM research setting:
\begin{equation}
\label{eq:phase0_setting}
\boxed{
\text{scalar output } p(t)\in\mathbb R,
\quad s=4,
\quad d_i=1,
\quad k=2,
\quad \text{natural boundary conditions only}.
}
\end{equation}
This setting will serve as the base model for:
\begin{enumerate}
  \item local existence and uniqueness,
  \item exact local linear-system construction,
  \item Hermite coefficient reconstruction,
  \item local-support proofs, and
  \item later comparison with global MINCO.
\end{enumerate}

\subsection{Why This Restriction Is Necessary}
\label{subsec:phase0_why}

The restriction in Phase~0 is deliberate rather than conservative.
The MINCO construction succeeds because it first fixes a unified multistage control-effort minimization problem, proves necessary and sufficient optimality conditions together with uniqueness, and only then introduces sparse parameterization and spatial--temporal deformation. Reproducing this discipline for BLOM requires first fixing one mathematically controlled regime. Trying to study arbitrary $s$, arbitrary $k$, multiple boundary conventions, and additional constraints from the beginning would mix several independent sources of difficulty and make later theorems ambiguous.

\subsection{Notation Table}
\label{subsec:phase0_notation}

\begin{table}[h]
\centering
\caption{Notation used in Phase~0.}
\label{tab:phase0_notation}
\begin{tabular}{ll}
\hline
Symbol & Meaning \\
\hline
$M$ & number of trajectory segments \\
$q_0,\dots,q_M$ & scalar waypoints \\
$T_1,\dots,T_M$ & positive segment durations \\
$\mathcal T$ & total duration, $\sum_{i=1}^{M} T_i$ \\
$p_i$ & polynomial on the $i$-th segment \\
$s$ & control order; fixed as $s=4$ in Phase~0 \\
$d_i$ & number of prescribed derivatives at intermediate time $t_i$ \\
$k$ & BLOM window width; fixed as $k=2$ in Phase~0 \\
$W(i,2)$ & local window associated with segment $i$ \\
$p^{(\ell)}$ & $\ell$-th derivative of $p$ \\
$\theta$ & sparse BLOM parameter vector \\
\hline
\end{tabular}
\end{table}

\subsection{Sparse Parameter Vector}
\label{subsec:phase0_theta}

Following the MINCO-style external parameterization, we define the Phase~0 sparse parameter vector as
\begin{equation}
\label{eq:phase0_theta}
\theta
=
(q_1,\dots,q_{M-1}, T_1,\dots,T_M)^\top.
\end{equation}
The boundary waypoints $q_0$ and $q_M$ are treated as fixed data.
In Phase~0, no optimization over $\theta$ is performed yet; instead, $\theta$ is only treated as the canonical input format for all subsequent theory and code modules.

\subsection{Canonical Extension Path}
\label{subsec:phase0_extension}

After the Phase~0 scalar theory is stabilized, the intended extension order is
\[
\mathbb R
\longrightarrow
\mathbb R^m,
\qquad
k=2
\longrightarrow
k\ge 2,
\qquad
\text{natural BC}
\longrightarrow
\text{alternative boundary mechanisms}.
\]
This order should not be reversed.

\paragraph{Phase~0 Deliverables.}
The mathematical deliverables of Phase~0 are:
\begin{enumerate}
  \item a formal problem statement,
  \item a notation table, and
  \item a canonical setting specification $(s,k,d_i)=(4,2,1)$ with scalar output and natural boundary conditions.
\end{enumerate}
